% readme_loop.tex -- the research loop, for README.md "## How a result is made".
%
% WHAT IT SHOWS
%   A clockwise racetrack with five stations: Run -> Verify -> Write up -> Index -> Plan -> Run.
%   Run sits OUTSIDE the dashed repository boundary (raw per-seed JSON lives on cluster scratch),
%   so a result re-enters the repository only through Verify. Write up produces dated findings;
%   an older finding keeps its full text under an outlined RETRACTED stamp and is superseded by
%   a newer card (additive, nothing deleted). A CI-tests chip gates Write up -> Index. Index is a
%   solid corpus JSON that renders a map and rebuilds a dashed, shadowless in-memory SQL store,
%   which the next Plan reads. In the centre, two peer agents stand on one working agreement
%   beside the maintainer, who owns scope and the issue tracker.
%
% RULES IT FOLLOWS (dd-readme.sty)
%   * achromatic: ddInk / ddMute / ddRule / ddWash / ddFixed / ddTrack and the shade greys only;
%     no red, no blue, no gold, no 1.6 pt element;
%   * sharp = data / state, rounded = the agents, grey fill = fixed machinery (the CI chip),
%     dashed = scope or ephemeral (the boundary, the in-memory store);
%   * 8 pt floor, drawn at final size, opaque white canvas, no result numbers, no counts;
%   * every glyph is a Type 1 outline: \fMono is re-pointed at txtt below, because the T1
%     \ttfamily of dd-readme.sty falls back to a bitmap (Type 3) font when cm-super is absent,
%     and pdftocairo then embeds those glyphs as PNG <image>s in the SVG;
%   * the canvas is 499 x 221 bp (176.0 x 78.0 mm), whole big points on purpose: pdftocairo
%     rescales a page whose size is not a whole number of bp (0.27 % at 176 x 78 mm), and that
%     shift alone fails check_render.py. Coordinates below are in mm.
%   * layer order: track -> chevrons -> boundary -> medallions (shadow over the band) ->
%     artifacts -> centre.
\documentclass[tikz,border=0pt]{standalone}
\usepackage{dd-readme}
% figure-local: an outline typewriter face (see RULES above); T1, so \{ \} are real braces
\renewcommand{\fMono}{\fontsize{8}{10}\fontfamily{txtt}\selectfont}

% ---- figure-local primitives ---------------------------------------------------------------
% track geometry: two straights (y = 18, 54) joined by half circles about (46,36) and (128,36)
\newcommand{\trackpath}{(46,18) -- (128,18) arc[start angle=-90, end angle=90, radius=18]
  -- (46,54) arc[start angle=90, end angle=270, radius=18] -- cycle}
% \stadium{r}: the same racetrack at offset radius r (the band is r = 16.7 .. 19.3)
\newcommand{\stadium}[1]{(46,36-#1) -- (128,36-#1) arc[start angle=-90, end angle=90, radius=#1]
  -- (46,36+#1) arc[start angle=90, end angle=270, radius=#1] -- cycle}
\newcommand{\medR}{6.5}
% station centres as numbers (a named coordinate ignores the shift inside \ddshadow)
\newcommand{\stationpts}{(28,36),(68,54),(108,54),(146,36),(88,18)}

% shadow steps for the part of a shadow that falls ON the track band: the band-colour
% equivalents of ddShadeFar / ddShadeNear (multiplied onto ddTrack), so the outer step reads
% as shade there too instead of as a light halo.
\definecolor{ddBandFar} {HTML}{D8D8D8}
\definecolor{ddBandNear}{HTML}{CBCBCB}
% \bandshadow{<path>}: \ddshadow, then the same two steps recoloured inside the band only
\newcommand{\bandshadow}[1]{%
  \ddshadow{#1}%
  \begin{scope}[even odd rule]
    \clip \stadium{19.3} \stadium{16.7};
    \begin{scope}[shift={(1.0mm,-1.0mm)}]\fill[ddBandFar]  #1;\end{scope}%
    \begin{scope}[shift={(0.5mm,-0.5mm)}]\fill[ddBandNear] #1;\end{scope}%
  \end{scope}}

% \chevron{(x,y)}{angle}: a filled Stealth-shaped arrowhead (3.0 x 2.4 mm) centred ON the
% track centreline, pointing along travel. A polygon, not an arrow tip, so no shaft shows.
\newcommand{\chevron}[2]{%
  \begin{scope}[shift={#1}, rotate=#2]
    \fill[ddInk] (1.5,0) -- (-1.5,1.2) -- (-0.6,0) -- (-1.5,-1.2) -- cycle;
  \end{scope}}

% \medallion{(x,y)}: shadowed white disc that carries a station icon
\newcommand{\medallion}[1]{%
  \bandshadow{#1 circle (\medR)}%
  \draw[ddInk, line width=0.75pt, fill=white] #1 circle (\medR);}

% \sheet{x}{y}: one 7 x 9 mm raw-result sheet with its south-west corner at (x,y)
\newcommand{\sheet}[2]{%
  \ddshadow{(#1,#2) rectangle ++(7,9)}%
  \draw[datum] (#1,#2) rectangle ++(7,9);}

% \textlines{x0}{x1}{y0}{dy}: three ddRule hairlines, the last one shorter
\newcommand{\textlines}[4]{%
  \draw[hair] (#1,#3) -- (#2,#3);
  \draw[hair] (#1,#3-#4) -- (#2,#3-#4);
  \draw[hair] (#1,#3-2*#4) -- ($(#1,#3-2*#4)!0.6!(#2,#3-2*#4)$);}

% \agentcard{sw-x}{sw-y}{name}: a rounded (learned-style) peer card with a '>_' prompt glyph.
% Its shadow is drawn separately (\agentshadow), before the slab, so the card stands ON it.
\newcommand{\agentshadow}[2]{\ddshadow{[rounded corners=1.2pt] (#1,#2) rectangle ++(22,11)}}
\newcommand{\agentcard}[3]{%
  \draw[learned] (#1,#2) rectangle ++(22,11);
  \draw[ddInk, line width=0.6pt, line cap=round, line join=round]
    (#1+3.0,#2+8.3) -- (#1+4.4,#2+7.3) -- (#1+3.0,#2+6.3)
    (#1+5.0,#2+6.3) -- (#1+7.0,#2+6.3);
  \node[font=\fLabel, text=ddInk] at (#1+11,#2+3.3) {#3};}

\tikzset{
  icon/.style   = {draw=ddInk, line width=0.6pt, line cap=round, line join=round},
  leader/.style = {draw=ddMute, line width=0.6pt, dotted},
  stsub/.style    = {font=\fSmall, text=ddMute},
  stname/.style   = {font=\fTitle, text=ddInk},
}

\begin{document}
\begin{tikzpicture}[x=1mm,y=1mm]
\sansmath
\ddcanvas{499bp}{221bp}

% ---- station anchors -----------------------------------------------------------------------
\coordinate (run)    at (28,36);
\coordinate (verify) at (68,54);
\coordinate (write)  at (108,54);
\coordinate (index)  at (146,36);
\coordinate (plan)   at (88,18);

% ==== track =================================================================================
\draw[ddTrack, line width=2.6mm, line cap=round] \trackpath;
\chevron{(88,54)}{0}
\chevron{(140.73,23.27)}{-135}
\chevron{(112,18)}{180}
\chevron{(64,18)}{180}
\chevron{(33.27,48.73)}{45}

% ==== repository boundary (crossed twice by the track, at x = 40) ===========================
\draw[scope] (40,1.5) rectangle (174,76.5);
% region names flank the boundary at its top-left corner, where the eye enters the figure,
% so "cluster | repository" is read before the off-repo fans are
\node[stsub, anchor=base east, inner sep=0pt] at (38.2,73.4) {cluster};
\node[stsub, anchor=base west, inner sep=0pt] at (41.8,73.4) {repository};

% ==== stations: medallions float over the track (shadow on top of the band) =================
\foreach \p in \stationpts {\medallion{\p}}

% RUN: server rack
\begin{scope}[shift={(run)}]
  \foreach \dy in {-2.2,0,2.2} {
    \draw[icon, rounded corners=0.5pt] (-3,\dy-0.8) rectangle (3,\dy+0.8);
    \fill[ddInk] (-2.0,\dy) circle (0.3);}
\end{scope}
\node[stname, anchor=east] at (19.5,37.8) {Run};
\node[stsub,  anchor=east] at (19.5,33.8) {multi-seed};

% VERIFY: shield with a tick
\begin{scope}[shift={(verify)}]
  \draw[icon] (-2.7,2.75) -- (2.7,2.75) -- (2.7,-0.25)
    .. controls (2.7,-1.3) and (1.4,-2.1) .. (0,-2.75)
    .. controls (-1.4,-2.1) and (-2.7,-1.3) .. (-2.7,-0.25) -- cycle;
  \draw[icon] (-1.3,0.3) -- (-0.3,-0.8) -- (1.4,1.3);
\end{scope}
\node[stname] at (68,66.5) {Verify};
\node[stsub]  at (68,62.3) {recompute or stop};

% WRITE UP: page with folded corner, text lines, pen stroke
\begin{scope}[shift={(write)}]
  \draw[icon] (-2.3,-3) -- (-2.3,3) -- (0.9,3) -- (2.3,1.6) -- (2.3,-3) -- cycle;
  \draw[icon] (0.9,3) -- (0.9,1.6) -- (2.3,1.6);
  \draw[ddInk, line width=0.4pt] (-1.4,1.2) -- (0.5,1.2) (-1.4,0) -- (1.4,0) (-1.4,-1.2) -- (0.6,-1.2);
  \draw[white, line width=2.2pt, line cap=round] (3.1,-0.6) -- (1.3,-2.9);   % halo: pen sits on the page
  \draw[ddInk, line width=0.8pt, line cap=round] (3.1,-0.6) -- (1.3,-2.9);
\end{scope}
\node[stname] at (108,66.5) {Write up};
\node[stsub]  at (108,62.3) {dated findings};

% INDEX: 3 x 3 table, header row filled
\begin{scope}[shift={(index)}]
  \fill[ddMute] (-3,0.8) rectangle (3,2.4);
  \draw[icon] (-3,-2.4) rectangle (3,2.4);
  \draw[icon] (-3,0.8) -- (3,0.8) (-3,-0.8) -- (3,-0.8) (-1,-2.4) -- (-1,2.4) (1,-2.4) -- (1,2.4);
\end{scope}
\node[stname, anchor=west] at (154.5,37.8) {Index};
\node[stsub,  anchor=west] at (154.5,33.8) {research map};

% PLAN: clipboard with two ticked lines
\begin{scope}[shift={(plan)}]
  \draw[icon] (-2.5,-3.2) rectangle (2.5,3.2);
  \draw[icon, fill=white] (-1.2,2.7) rectangle (1.2,3.7);
  \draw[icon] (-1.7,1.1) -- (-1.2,0.6) -- (-0.5,1.4);   \draw[icon] (0.2,1.0) -- (1.6,1.0);
  \draw[icon] (-1.7,-1.1) -- (-1.2,-1.6) -- (-0.5,-0.8); \draw[icon] (0.2,-1.2) -- (1.6,-1.2);
\end{scope}
\node[stname] at (88,8.9) {Plan};
\node[stsub]  at (88,4.6) {next experiment};

% ==== RUN artifacts: arm and control fans, one job ==========================================
\foreach \x in {5,19} {
  \foreach \k in {2,1,0} {\sheet{\x+\k}{56+\k}}
  \node[font=\fMono, text=ddInk] at (\x+3.5,60.5) {\{\,\}};}
\node[font=\fSmall, text=ddInk, anchor=base] at (8.5,52.4)  {arm};
\node[font=\fSmall, text=ddInk, anchor=base] at (22.5,52.4) {control};
\draw[ddMute, line width=0.6pt, decorate, decoration={brace, amplitude=1.2mm}]
  (5,68.5) -- (28,68.5);
\node[font=\fSmall, text=ddInk] at (16.5,71.8) {same job};
% leader: Run medallion (0.5 mm off the ring) -> the gap between the two captions
\draw[leader] (23.59,41.44) -- (15.6,53.2);

% ==== WRITE UP artifacts: retracted older card, newer card, supersedes =====================
% older card: wide enough that the stamp keeps a clear margin on both sides, and four text
% lines so the text visibly runs on around the stamp (kept, not deleted)
\ddshadow{(142,60.5) rectangle (170,71.5)}
\draw[datum] (142,60.5) rectangle (170,71.5);
\draw[hair] (148,70.1) -- (167.5,70.1) (148,67.7) -- (167.5,67.7) (148,65.3) -- (167.5,65.3)
  (148,62.9) -- (156.8,62.9);
\node[rotate=-12, draw=ddMute, line width=0.6pt, fill=white, inner sep=0.8mm,
      font=\fSmall\bfseries, text=ddMute] at (158.75,65.4) {RETRACTED};
% newer card, drawn after the stamp
\ddshadow{(130,58) rectangle (146,69)}
\draw[datum] (130,58) rectangle (146,69);
\textlines{132.5}{143.5}{66.4}{2.6}
% supersedes: a permanent, additive link, so solid (dashed means scope or ephemeral here)
\draw[flowS] (136,69) |- (142,70.5);
\node[stsub, anchor=east] at (133.5,70.5) {supersedes};
\draw[leader] (114.5,56.5) -- (130,60);

% ==== CI gate on the Write up -> Index arc: fixed machinery, so grey-filled ================
\bandshadow{(134.73,46.23) rectangle ++(12,5)}
\node[fixedbox, minimum width=12mm, minimum height=5mm, inner sep=0pt, font=\fSmall, text=ddInk]
  at (140.73,48.73) {CI tests};

% ==== INDEX artifacts: solid corpus -> map -> dashed in-memory SQL -> the next Plan ==========
% the three labels share one baseline (y 6.2)
\ddshadow{(147,3.5) rectangle (166,10.5)}
\draw[datum] (147,3.5) rectangle (166,10.5);
\node[font=\fSmall, text=ddInk, anchor=base] at (156.5,6.2) {corpus JSON};
\ddshadow{(128.5,3.5) rectangle (142.5,10.5)}
\draw[datum] (128.5,3.5) rectangle (142.5,10.5);
\node[font=\fSmall, text=ddInk, anchor=base] at (135.5,6.2) {map .md};
% the cylinder keeps its shape from a phantom; the label is a separate node set on the
% shared baseline, clear of both end arcs
\node[cylinder, shape border rotate=90, aspect=0.2, minimum width=24mm, minimum height=9mm,
      inner sep=0pt, draw=ddMute, line width=0.6pt, dash pattern=on 2pt off 1.5pt, fill=white,
      font=\fSmall] (sql) at (112.5,7.5) {\phantom{SQL, in memory}};
\node[font=\fSmall, text=ddInk, anchor=base] at (112.5,6.2) {SQL, in memory};
\draw[flowS] (147,7) -- (142.5,7);
\draw[flowS] (128.5,7) -- (124.7,7);
% the hand-off: the next Plan reads the store (upper-left shoulder of the drum -> Plan rim)
\draw[flowS] (100.6,11.5) -- (94.2,14.8);
\draw[leader] (147.97,29.28) -- (153.5,10.5);
\node[font=\fMono, text=ddInk] at (116,21.8) {settled?\ \ superseded?};

% ==== CENTRE: peer agents standing on one agreement, and the maintainer =====================
\ddshadow{(54,28.6) rectangle (102,33)}
\agentshadow{55}{33}
\agentshadow{79}{33}
\draw[sharp corners, draw=ddInk, line width=0.6pt, fill=ddWash] (54,28.6) rectangle (102,33);
\node[font=\fSmall, text=ddInk, anchor=base] at (78,30.1) {one working agreement};
\agentcard{55}{33}{Claude Code}
\agentcard{79}{33}{Codex}

\draw[ddInk, line width=0.6pt, fill=white]
  (111.5,35) .. controls (111.5,39) and (120.5,39) .. (120.5,35) -- cycle;
\draw[ddInk, line width=0.6pt, fill=white] (116,41.5) circle (2.0);
% 'maintainer' shares the slab label's baseline
\node[font=\fLabel, text=ddInk, anchor=base] at (116,30.1) {maintainer};
\node[stsub, anchor=base] at (116,26.6) {scope\,$\cdot$\,issues};

\end{tikzpicture}
\end{document}
